The relative Virasoro vacuum complex has classes in degrees zero and two.
The full unital Virasoro chiral algebra has chiral-chain classes on the projective line in degrees zero, one, three, and four.
A comparison map must therefore lose the relative degree-two class.
An explicit one-point chain gives such a map and its null-homotopy.

\section{The relative complex and its absolute primitive}

Work over $\mathbb C$, using algebraic tensor products and direct sums.
All complexes are cohomological, with $K[r]^j=K^{j+r}$ and differential $(-1)^r d_K$.
Ghost degree is cohomological degree.
Conformal weight is a separate integer grading.
Interchanging homogeneous factors of degrees $a,b$ contributes $(-1)^{ab}$.

Let $A=\operatorname{Vir}_{26}\otimes bc$ use the universal matter vacuum module and the conformal ghost vacuum $\Omega$.
The ghost relations and annihilators are
\[
 \{b_m,c_n\}=\delta_{m+n,0},\qquad
 \{b_m,b_n\}=\{c_m,c_n\}=0,
 \qquad b_n\Omega=0\ (n\geq-1),\quad c_n\Omega=0\ (n\geq2).
\]
The matter vacuum is implicit in each displayed ghost state.
The generators $b_n,c_n$ have degrees $-1,1$ and weights $-n$.
Set
\[
 Q=\sum_n c_{-n}L_n
 -\frac12\sum_{m,n}(m-n):c_{-m}c_{-n}b_{m+n}:,
 \qquad H=\mathcal L_0=L_0+L_0^{\mathrm{gh}}.
\]
Normal ordering moves creators to the left and omits contractions.
The sums are locally finite on algebraic states.
Theorem~\ref{thm:vir-square} and Lemma~\ref{lem:vir-Q-contractions} give
$Q^2=0$ and $\{Q,b_0\}=H$.

Fix a coordinate $z$ and define the vector-space complex
\[
 R=\ker b_0\cap\ker H\subset A.
\]
The weight-zero absolute complex has basis
\[
 e=\Omega,\qquad u=c_0\Omega,\qquad
 v=c_{-1}c_1\Omega,\qquad w=c_{-1}c_0c_1\Omega.
\]
Its differential and relative subcomplex are
\begin{equation}\label{eq:rp-four}
 Qe=Qv=Qw=0,\qquad Qu=-2v,
 \qquad (R,Q)=\mathbb Ce\oplus\mathbb Cv.
\end{equation}
The degrees of $e,u,v,w$ are $0,1,2,3$.
Indeed, matter weights are nonnegative, matter weight one vanishes, and the only negative ghost weight is that of $c_1$.
These facts leave exactly the four displayed weight-zero monomials.
The operator $b_0$ removes the terms containing $c_0$.
Its anticommutator with $Q$ contracts every nonzero absolute weight.
Consequently
\[
 H^\bullet(R)=\mathbb C[e]\oplus\mathbb C[v],\qquad
 H^\bullet(A,Q)=\mathbb C[e]\oplus\mathbb C[w].
\]
The formula $v=Q(-u/2)$ determines the relative class image under any comparison that extends to absolute states.

\section{A normalized chain in the one-point summand}

Set $X=\mathbb P^1$, with ordered cover
$U_0=\operatorname{Spec}\mathbb C[z]$, $U_\infty=\operatorname{Spec}\mathbb C[z^{-1}]$.
Their intersection is $U_{01}=\operatorname{Spec}\mathbb C[z,z^{-1}]$.
Let $\mathcal A^\ell$ be the left state module with connection $\partial_z+\partial$ and let $\mathcal A=\omega_X\otimes\mathcal A^\ell$.
Here $\partial$ is vertex translation.
Coordinate descent and the differential are specified in Section~\ref{sec:cbc26-unital-p1}.

Use the two resolutions $P,\mathcal Q$ from the construction of~\eqref{eq:uc-chain}.
Explicitly,
\[
 P^0=\operatorname{Sym}^{>0}_{\mathcal O_X}(\mathcal D_X),\qquad
 P^{-1}=\ker(P^0\longrightarrow\mathcal O_X),
 \qquad d_P:P^{-1}\hookrightarrow P^0.
\]
The augmentation on a generating differential operator $D$ is $D(1)$.
Write $p=[1_{\mathcal D}]\in P^0$ for the global generator given by the identity differential operator.
Thus $d_Pp=0$ and the augmentation sends $p$ to $1$.
The symbol $p$ denotes a degree-zero element of the nonunital resolution.

The degree-one term of the two-chart Thom--Sullivan resolution is
\[
 \mathcal Q^1=j_{01*}\mathcal O_{U_{01}}[t]\,dt.
\]
In degree zero it consists of triples $(f_0,f_\infty,F(t))$ with $F(0)=f_0$ and $F(1)=f_\infty$ on the overlap.
The differential differentiates $F(t)$.
Its connection differentiates the coefficients and fixes $t,dt$.
The variable $t$ here is an auxiliary interval coordinate.
Normalize $\int_0^1dt=1$ and $\operatorname{Res}_{z=0}(dz/z)=1$.

For a right $\mathcal D$-module $N$, put $h(N)=N\otimes_{\mathcal D_X}\mathcal O_X$, degreewise.
The one-point subcomplex of the actual chiral chains is
\[
 S(\mathcal A)=\Gamma(X,h(\mathcal A\otimes P\otimes\mathcal Q))[1]
 \lhook\joinrel\longrightarrow\mathsf C_X(\mathcal A).
\]
The target is the Chevalley--Cousin complex~\eqref{eq:uc-chain}, with its full positive-arity unit carrier.
The one-point inclusion is a chain map.
A collision merges two distinct blocks, so there is no outgoing collision from a one-block term.
This is also immediate from Beilinson--Drinfeld~\cite{BD026}, equations (4.2.12.2) and (4.2.12.5), pp.~306--307.

Write $\varsigma$ for suspension, so $|\varsigma x|=|x|-1$ and $d(\varsigma x)=-\varsigma(dx)$.
The section $(dz/z)\otimes p\otimes dt$ is global in the term supported on $U_{01}$.
Define
\begin{equation}\label{eq:rp-tau}
 \tau=\varsigma\bigl[(dz/z)\otimes p\otimes dt\bigr]
 \in S(\omega_X)^0.
\end{equation}
Square brackets in this formula denote the image in $h$.

\begin{lemma}\label{lem:rp-unit}
The element $\tau$ is closed.
Its image in $H^0\mathsf C_X(\omega_X)$ has trace $1$ and is the chiral unit class.
\end{lemma}

\begin{proof}
Both $p$ and $dt$ are closed.
The one-point term has no outgoing collision.
These observations prove closedness directly in the actual complex.

The right de Rham convention is
$\operatorname{DR}_X(\omega_X)=\Omega_X^\bullet[1]$.
Consequently $S(\omega_X)$ represents $R\Gamma(X,\Omega_X^\bullet)[2]$.
To identify the specific cycle, use the intermediate complex
\[
 \Gamma\bigl(X,\operatorname{DR}_X(\omega_X\otimes P\otimes\mathcal Q)\bigr)[1].
\]
The same tensor defines a closed element in its terminal Spencer term.
The Spencer augmentation maps it to~\eqref{eq:rp-tau}.
The augmentation of $P$ maps it to the form $(dz/z)\,dt$.
The two maps are the resolution comparisons used in Beilinson--Drinfeld, \S4.2.12, p.~306.
For the right action on $\omega_X$, the Spencer differential sends $f\,dz\otimes\partial_z$ to $-f'\,dz$.
The identifications $f\,dz\otimes\partial_z\mapsto f$ and $f\,dz\mapsto f\,dz$ give $\operatorname{DR}_X(\omega_X)=\Omega_X^\bullet[1]$.
In particular, the top form has its displayed positive sign.
Integrating the interval form gives the \v Cech one-cocycle $dz/z$ with values in $\Omega_X^1$.
It has trace $1$ under the stated residue normalization.
The trace of the singleton term agrees with the trace on the Cousin diagram by Beilinson--Drinfeld, \S4.3.3, pp.~315--316.
That proposition identifies $H^0\mathsf C_X(\omega_X)$ with $\mathbb C$ and all other unit cohomology with zero.
Thus the image of $\tau$ is the normalized unit class.
\end{proof}

\section{The actual comparison and its class images}

For an algebraic state $a\in A$, let $a_z$ be the constant section in the $z$-coordinate trivialization over $U_0$.
Restriction makes it a section over $U_{01}$.
For homogeneous $a$, define
\begin{equation}\label{eq:rp-J}
 J_z(a)=(-1)^{|a|}\varsigma
 \bigl[(dz/z)\otimes a_z\otimes p\otimes dt\bigr]
 \in S(\mathcal A)^{|a|}\subset\mathsf C_X(\mathcal A)^{|a|}.
\end{equation}
Extend by linearity to all algebraic states.
All tensors in this formula are finite.
The sole pole $z^{-1}$ is regular on the specified overlap.
No extension of $a_z$ across infinity is required by $\mathcal Q^1$.
For the Poincar\'e--Birkhoff--Witt (PBW) filtration, give matter and $b$ creators order one and $c$ creators order zero.
The auxiliary resolutions have order zero.

\begin{proposition}\label{prop:rp-J}
The map $J_z:(A,Q)\to\mathsf C_X(\mathcal A)$ is a degree-preserving map of vector-space complexes.
It preserves the PBW filtration from~\eqref{eq:uc-symbolmap}.
Its restriction $\Phi_z=J_z|_R$ satisfies
\begin{equation}\label{eq:rp-null}
 \Phi_z(e)=1_{\mathcal A*}(\tau),\qquad
 \Phi_z(v)=d\bigl(-J_z(u)/2\bigr).
\end{equation}
Here $1_{\mathcal A*}$ is the actual chain map induced by the chiral unit $\omega_X\to\mathcal A$.
\end{proposition}

\begin{proof}
In the unsuspended tensor, only the differential of $a$ contributes.
The remaining factors have $d_Pp=0$ and $d_{\mathcal Q}dt=0$.
For $|a|=k$, suspension therefore gives
\[
 dJ_z(a)=(-1)^{k+1}\varsigma
 \bigl[(dz/z)\otimes(Qa)_z\otimes p\otimes dt\bigr]
 =J_z(Qa).
\]
Passing to $h$ preserves this equation because $Q$ commutes with the connection.
There is no outgoing collision from its one-point image.
This proves the chain-map identity in the full complex.
The auxiliary factors have PBW order zero, so $J_z$ preserves the state order.
The first formula in~\eqref{eq:rp-null} follows by substituting $a=e$.
The second follows from $Qu=-2v$ in~\eqref{eq:rp-four}.
\end{proof}

The null-homotopy is also explicit as a map of complexes.
Let $\Phi_z^0$ send $e$ to $1_{\mathcal A*}(\tau)$ and $v$ to zero.
Define the degree-minus-one linear map $K:R\to\mathsf C_X(\mathcal A)$ by
\[
 K(e)=0,\qquad K(v)=-J_z(u)/2.
\]
Then $\Phi_z-\Phi_z^0=dK+KQ$.
These are statements about vector-space complexes and do not assert multiplicativity of $J_z$.

\begin{theorem}\label{thm:rp-images}
Use the degree and trace conventions above.
Identify the target cohomology by Theorem~\ref{thm:uc-global} as
$\Lambda(\xi_1,\xi_3)$, with degrees $1,3$.
Then the actual relative comparison~\eqref{eq:rp-J} induces
\[
 H(\Phi_z)[e]=1,\qquad H(\Phi_z)[v]=0.
\]
Its graded kernel, image, and cokernel are
\begin{align*}
 \ker H(\Phi_z)&=\mathbb C[v], & |[v]|&=2,\\
 \operatorname{im}H(\Phi_z)&=\mathbb C1,\\
 \operatorname{coker}H(\Phi_z)&=
 \mathbb C\xi_1\oplus\mathbb C\xi_3\oplus\mathbb C\xi_1\xi_3.
\end{align*}
The cokernel degrees are $1,3,4$.
The same assertions hold for the induced map to PBW associated graded chiral chains.

For the absolute comparison $J_z$, normalize $\xi_3$ by the image of $\tau\otimes\eta$, where $\eta=\alpha\beta\gamma$.
Then $H(J_z)[w]=\xi_3$.
Thus $H(J_z)$ is injective, with image $\mathbb C1\oplus\mathbb C\xi_3$ and cokernel $\mathbb C\xi_1\oplus\mathbb C\xi_1\xi_3$.
\end{theorem}

\begin{proof}
The chain identities~\eqref{eq:rp-null} give the relative images.
Lemma~\ref{lem:rp-unit} and the unital zigzag~\eqref{eq:uc-zigzag} show that the vacuum image is nonzero.
The source has precisely the two classes in~\eqref{eq:rp-four}.
The target basis in Theorem~\ref{thm:uc-global} therefore gives the stated kernel and cokernel.

For the absolute degree-three class, all four weight-zero states lie in the ghost subalgebra $B$.
The one-point construction commutes with the inclusion $B\hookrightarrow A$ and the global evaluation $\epsilon:B\to E_X$.
Here
\[
 E=\Lambda(\alpha,\beta,\gamma),\quad
 d\alpha=-\alpha\beta,\quad d\beta=-2\alpha\gamma,\quad d\gamma=-\beta\gamma,
 \qquad f(z)=\alpha+z\beta+z^2\gamma.
\]
The identification $S(E_X)=S(\omega_X)\otimes E$ moves $E$ past $P\otimes\mathcal Q$ with the Koszul sign.
The evaluation formulas of~\eqref{eq:uc-eval} give
\begin{equation}\label{eq:rp-Eimages}
 \epsilon(e)=1,\qquad
 \epsilon(u)=-\beta-2z\gamma,\qquad
 \epsilon(v)=-\alpha\gamma-z\beta\gamma,\qquad
 \epsilon(w)=\eta.
\end{equation}
In particular, $\epsilon(v)=d((\beta+2z\gamma)/2)$ at every $z$.
The degree-three image is constant on both charts.
After moving the coefficient algebra past $dt$, the Koszul sign is $(-1)^3$.
It cancels the prefactor in~\eqref{eq:rp-J}.
Hence $\mathsf C_X(\epsilon)J_z(w)=\tau\otimes\eta$ in the constant coefficient one-point term.

The generator complex for $\omega_X[-3]$ represents $R\Gamma(X,\Omega_X^\bullet)[-1]$.
Its top de Rham class is $\tau\otimes\eta$, of degree three.
The positive symmetric-degree comparison~\eqref{eq:uc-free-map} sends this class to $\xi_3$.
It is nonzero by that quasi-isomorphism.
This determines both absolute class images and their kernel and cokernel.

Finally, $e,u,v,w$ have PBW order zero.
The symbol of~\eqref{eq:rp-J} is the same tensor formula with the state replaced by its symbol.
The diagram with~\eqref{eq:uc-symbolmap} therefore commutes on actual chains.
The PBW target cohomology lies entirely in order zero, and its unital comparison has the same normalized classes.
This proves the associated graded assertion.
\end{proof}

\section{Coordinate dependence and the relative descent obstruction}

The map $J_z$ uses an ordered affine cover, its interval resolution, and the coordinate trivialization on $U_0$.
The cycle $\tau$ is a chosen algebraic representative of the top de Rham class.
These choices define a map from the fixed-coordinate state complex to the unmarked global chiral-chain complex.
The construction does not require a marked-point coefficient module in the target.

A change of coordinate does not preserve the fixed relative subspace by restriction of the ordinary state transitions.
For the convention $z\mapsto z/(1+sz)$ acting by $g_s=\exp(s\mathcal L_1)$,
\[
 \mathcal L_1v=-c_0c_1\Omega,\qquad
 \mathcal L_1(c_0c_1\Omega)=0.
\]
The ghost commutator $[\mathcal L_m,c_n]=(-2m-n)c_{m+n}$ gives both identities.
Consequently
\begin{equation}\label{eq:rp-transition}
 g_sv=v-sc_0c_1\Omega,\qquad
 b_0g_sv=-sc_1\Omega,\qquad Hg_sv=sc_0c_1\Omega.
\end{equation}
For $s\ne0$, this vector is outside $R$.
Thus these transition operators do not define a relative subbundle with model fiber $R$.

Any map of vector-space complexes from $R$ to this unmarked target must kill its degree-two cohomology class.
The target degree-two cohomology vanishes.
A map sending $[e]$ to the normalized unit has the class images in Theorem~\ref{thm:rp-images}, regardless of its chain representative.
Equation~\eqref{eq:rp-null} supplies a primitive for the specific comparison constructed here.
Neither this degree obstruction nor~\eqref{eq:rp-transition} excludes a different marked or geometric relative complex.
Such a complex requires defined coefficient sheaves, transported constraints, and compatible transition maps before a global comparison can be stated.
